\section{Experimento 1}

Para este primer experimento buscaremos obtener una forma de onda cuadrada para analizar.\\
Para ello primero calibramos la punta de prueba y luego conectamos la salida del generador de funciones a la entrada del osciloscopio.\\
Pondremos las siguientes configuraciones en el generador de funciones: frecuencia de 1 kHz, sin CC y con media vuelta de perilla del nivel de salida.\\
Ahora verificamos que el osciloscopio tenga la siguiente configuración: acople en CC, sin límite de ancho de banda, sonda en x10 y sin inversión de la señal.\\

IMAGEN\\

Ahora pondremos el menú matemático del osciloscopio y seleccionaremos la opción de FFT, para el CH1 con zoom en x1.\\

Ahora veremos el espectro de la señal de entrada, y con el control de Posición Horizontal, pondremos el espectro en el centro de pantalla para luego poner el zoom en x10 y así poder ver mejor la señal.\\

OBSERVACIONES\\

Ahora con los Cursores propios del osciloscopio para frecuencia, mediremos las componentes espectrales de la señal, para ello seleccionaremos los siguientes ajustes:

\begin{itemize}
    \item Modo de adquisición: Promedio
    \item Número de cuentas: 64
    \item Ventana: Rectangular
\end{itemize}

CUADRO\\

Ahora cambiamos la ventana a Flattop y con los cursores de amplitud mediremos las amplitudes de las 7 primeras componentes espectrales de la señal, al estar en dBv, para poder calcular la amplitud en V, utilizaremos la siguiente ecuación:

\begin{equation*}
    V = 10^{\frac{dBv}{20}}
\end{equation*}

CUADRO\\

Y con estos datos calculamos el valor eficaz:

\begin{equation*}
    V_{ef} = \sqrt{\sum_{i=1}^{n} V_i^2}
\end{equation*}

